\section{Evaluation}

After the modeling stage reported all the necessary performance metrics, the evaluation stage discusses the compliance of the results with the project success criteria. The assessment plan is organized as follows: in the \textit{Evaluate Results} subsection, the key performance criteria are evaluated both in terms of data mining and business success criteria; in the \textit{Review Process} the quality of the data mining process is reported; and in the \textit{Determine Next Steps} we provide the final conclusions about the further project development.

\subsection{Evaluate Results}

\subsubsection{Assessment of Data Mining Results w.r.t. Business Success Criteria}

\begin{itemize}
  \item \textbf{Conversion rate prediction:} The Histogram Gradient Boosting Regressor demonstrated a sufficiently smaller MAE on the test set (0.1011) compared to the baseline approach (0.1551), as well as a positive $R^2$ score of 0.4333. In business terms, the regression model achieves \textbf{reliable} repository-level fork-to-PR conversion metric computed from public event data.
  \item \textbf{Classification:} The best classification model was Histogram Gradient Boosting as well, the model exceeded the threshold of a 5\% absolute increase in prediction performance compared to the baseline rate (macro F1 increase of 44\%). From the business point of view, the predictive model identifies high- and low-conversion repositories \textbf{with a sufficiently better} accuracy than a naive baseline.
  \item \textbf{Ranking quality:} The ranking quality at 10\% of the classification model appears to be sufficient, exceeding the random ranking rate precision of 47.24\% by more than 45 percentage points. The lift is almost 2 accordingly, and the near‑perfect NDCG@10\% (0.952) confirms that the model places converting repositories overwhelmingly at the beginning of the ranked list. Regarding the business objectives, this performance provides an evidence towards \textbf{reliability} of the conversion predictions with respect to the ranking produced by the model.
  \item \textbf{Validation stability:} The regression and classification models were evaluated on two temporal folds (training on winter-spring and evaluation on summer-autumn periods and vice versa). For macro-F1, precision@10\%, NDCG@10\%, and MAE, the relative degradation between the best and worst fold remained below 15\% (3.8\%, 3.3\%, 3.0\%, and 5.7\%, respectively), satisfying the stability criterion. However, the \(R^2\) score decreased by 28.5\% (from 0.372 to 0.266), exceeding the threshold of 15\%.

  For the regression task, the success criterion requires only a positive \(R^2\) on the test fold, not a degradation limit. The model achieves positive \(R^2\) values in both folds , and the absolute prediction error (MAE) remains stable. Therefore, the regression model is still \textbf{temporally robust} for its primary purpose, despite the relative drop in the explained variance. For classification, macro-F1 stability is well within the acceptable range, confirming that the model’s ability to identify high- and low-conversion repositories is \textbf{consistent across seasons}. Moreover, the ranking performance remains mostly on the same level, which also leads to consistency.
  \item \textbf{Interpretability:} The feature importance analysis (permutation importance) reveals that for both the classification and regression models, the strongest driver of predicted conversion is the number of issue comments. In the classification model, shuffling this feature reduces F1 score by 0.27, while all other features (watches, pushes, issues, season) contribute less than 0.03, i.e., an order of magnitude smaller. For the regression model, the same feature increases MAE by 0.053 when permuted; however, watches and pushes also show meaningful impact (0.039 and 0.012, respectively, corresponding to 73\% and 23\% of the top effect), while issues and season remain negligible.

  The performance of both models are also verified in terms of SHAP values. The repositories samples of the test set were analyzed in terms of the impact of the characteristics of each model, and the results are reported in Figure \ref{fig:shap_cls} and Figure \ref{fig:shap_reg}.

    \begin{figure}
        \centering
        \includegraphics[width=0.8\linewidth]{shap_waterfall_cls.png}
        \caption{SHAP Waterfall for Classification}
        \label{fig:shap_cls}
    \end{figure}

    \begin{figure}
        \centering
        \includegraphics[width=0.8\linewidth]{shap_waterfall_reg.png}
        \caption{SHAP Waterfall for Regression}
        \label{fig:shap_reg}
    \end{figure}
  
  From a business perspective, \textbf{repositories with active issue discussions} are the most likely to turn forks into pull requests. For classification (high/low conversion), no other feature matters meaningfully. For estimating the exact conversion rate (regression), watch counts (popularity) and push activity (code churn) also provide non‑negligible signals. Still, issue comments dominate, and the overall explanation remains interpretable: engaged discussions signal a healthy, welcoming community that encourages contributions. This satisfies the business success criterion.
  \item \textbf{Sensitivity check:} The conversion rate rankings computed with 30‑day and 90‑day observation windows are almost perfectly correlated (Spearman = 0.991, Kendall = 0.980). No repository changes from the top to the bottom decile between the two windows. Only 3.6\% of conversions occur exclusively between day 31 and day 90, and the median conversion rate does not shift (0.0). From a business perspective, this fact means that the label definition (fork‑to‑PR matching within a window) is robust: \textbf{the heuristic produces a consistent and reproducible repository‑level conversion metric}, and the binary classification target (has conversion) is \textbf{stable across plausible window lengths}. Stakeholders can trust that the ranking of repositories by conversion potential does not depend on an arbitrary choice of observation period, satisfying the sensitivity success criterion.
\end{itemize}

The main business success criteria were successfully covered by the modeling stage. Regarding pipeline documentation, all code, SQL queries and configuration parameters are versioned and publicly available in the GitHub project organization. Each stage (data extraction, feature engineering, labeling, modeling, evaluation) is implemented as a separate, self‑contained notebook with inline comments that explain how to rerun the pipeline on new time periods. This organization satisfies the business success criterion of a clearly documented and reproducible pipeline.

\subsubsection{Addressing Business Questions}
The main business objective, determining the fraction of forks that lead to a pull request and identifying which repository characteristics predict high or low conversion, is fully met. The regression model reliably estimates the conversion fraction, and the classification model shows that issue comment activity is the dominant distinguishing feature. Concerning the fraction of forks that lead to at least one PR, the regression metric is reliable and stable across observation windows.  
Regarding the characteristics that distinguish converting repositories, issue comment activity overwhelmingly drives high conversion; for classification, no other feature adds meaningful signal.
As for the existence of repository segments with consistently low conversion despite high fork activity, the total estimated number calculated out of 10,000 repositories is 9,158 (9.16\%). A repository was considered high‑fork if its fork count ($\texttt{n\_forks}$) was at or above the 75th percentile, and low‑conversion if its conversion rate ($\texttt{fork\_conversion\_rate}$) was at or below the 25th percentile. In many cases, such repositories are educational projects or project configurations. High fork activity alone does not signal a healthy and contribution‑ready community.

\subsubsection{Approved Models}
The business objectives require two distinct predictive models: repository-level fork-to-PR conversion metric regression and high- and low-conversion repositories classification. Based on this evaluation, the Histogram Gradient Boosting is suggested for both the regression objective and the classification task.

\subsection{Review Process}

The data mining process was reviewed for correctness, adherence to assumptions, reproducibility, and completeness with respect to the business objectives.

\begin{itemize}
    \item \textbf{Data extraction and sampling:} The multi‑season sampling (winter‑spring, summer‑autumn) and repository filters (minimum 5 forks, at least one interaction signal) were applied consistently. The random sample of 100,000 repository–season pairs is representative of active repositories. No data leakage occurred between feature extraction and label construction. This ensures that business decisions based on the model’s outputs are grounded in a representative and unbiased sample of the repository population.

    \item \textbf{Fork‑to‑PR matching heuristic:} The rule (same user, timestamp after fork, within 30/90 days) is clearly defined and was shown to be stable across observation windows (Spearman rank correlation 0.99). However, the heuristic was not quantitatively validated against a manually labelled sample. While this does not affect the relative ranking of repositories (the primary business need), it may bias absolute conversion fraction estimates. For the business goal of prioritising community investment, the ranking stability is sufficient; for precise conversion estimates, a manual validation would be required in a future iteration.

    \item \textbf{Temporal splits and stability:} Chronological train/test splits (first two months train, third month test) and two‑fold seasonal cross‑validation respect time order, avoiding future information leakage. Stability across seasons was confirmed for all key metrics except $R^2$ (which is not required to be stable). This means the business can rely on the model’s ranking and classification performance across different periods without frequent retraining.

    \item \textbf{Class imbalance and data preparation choices:} The positive conversion rate in the sample is ~48\% (due to the minimum fork filter), so no special balancing was needed. Winsorization (99th percentile) and log1p transformation were applied correctly to handle heavy tails, preserving the full sample while stabilising variance. These choices were sufficient to enable robust modeling.

    \item \textbf{Modeling approach:} Histogram Gradient Boosting was used for both regression and classification, which is appropriate for tabular, non‑linear, skewed data. Baseline comparisons (mean prediction for regression, stratified random for classification) confirm meaningful improvement. The modeling approach directly supports the business need to predict conversion rates and classify high/low conversion repositories.

    \item \textbf{Evaluation method:} The evaluation covered all data mining success criteria: regression (MAE, $R^2$), classification (macro‑F1, baseline improvement), ranking (precision@10\%, NDCG, lift), temporal stability, interpretability (permutation importance, SHAP values), and sensitivity (30d vs 90d).

    \item \textbf{Assumptions and data problems:} The key assumption that fork‑PR matching via user identity is reliable has not been empirically validated, but consistency checks support its use for ranking. No other misleading assumptions were identified. The only data problem was the large volume, which was mitigated by early sampling.

    \item \textbf{Improvements for a future iteration:} If the project were repeated, we would add a small manual verification of the matching heuristic (e.g., 100 random cases) to quantify false‑positive / ‑negative rates, include the segmentation analysis for high‑fork low‑conversion repositories, collect additional repository metadata (e.g., language, age) to potentially improve regression $R^2$ stability, train new models (SVM, neural networks), and fully automate the pipeline for quarterly refreshes.
\end{itemize}

\subsection{Determine Next Steps}

Based on the evaluation, the project meets all the business success criteria, and the model’s ranking performance is useful for prioritization purposes. However, there is scope for improvement: the fork-to-PR matching heuristic could undergo manual validation, and additional repository features (e.g. primary language and age) could be collected to potentially enhance regression stability. No additional CRISP-DM iteration is required at this stage, as the current results satisfy the stated goals and the project can proceed to deployment. While more data is not necessary, as the sampling strategy has already produced a representative set, a small manually labeled sample could provide additional evidence of the heuristic’s accuracy. The model is ready for deployment and can be used to rank repositories by conversion potential and flag high-fork, low-conversion segments. The main recommendation is to test the model in practice by selecting the top 10\% of predicted converting repositories, tracking their actual pull-request activity over several months and comparing it with the model’s predictions. This real-world feedback will confirm the business value and guide any future refinements.

